\documentclass[12pt,a4paper]{article}
\usepackage[utf-8]{inputenc}
\usepackage[margin=1in]{geometry}
\usepackage{graphicx}
\usepackage{amsmath}
\usepackage{amssymb}
\usepackage{array}
\usepackage{tabularx}
\usepackage{xcolor}
\usepackage{hyperref}
\usepackage{fancyhdr}
\usepackage{booktabs}

\pagestyle{fancy}
\fancyhf{}
\rhead{Digi Herba - Database Design}
\lhead{ERD, LRS \& Transformasi}
\cfoot{\thepage}

\title{\textbf{Desain Database Digi Herba}\\Entity-Relationship Diagram, Logical Record Structure, \& Transformasi}
\author{Backend Architecture Documentation}
\date{\today}

\begin{document}

\maketitle
\newpage

\tableofcontents
\newpage

% ============================================================================
\section{Entity-Relationship Diagram (ERD)}
% ============================================================================

\subsection{Deskripsi Entitas Utama}

Berikut adalah entitas-entitas utama dalam sistem Digi Herba:

\subsubsection{Entitas Core}
\begin{enumerate}
    \item \textbf{Users} - Menyimpan data pengguna aplikasi
    \item \textbf{Sessions} - Menyimpan session pengguna yang aktif
    \item \textbf{Password Reset Tokens} - Token untuk reset password
    \item \textbf{User Providers} - Integrasi social login
\end{enumerate}

\subsubsection{Entitas Otorisasi \& Permissions}
\begin{enumerate}
    \item \textbf{Roles} - Peran pengguna (admin, user, moderator, dll)
    \item \textbf{Permissions} - Izin/akses spesifik
    \item \textbf{Model Has Roles} - Relasi many-to-many Users \& Roles
    \item \textbf{Model Has Permissions} - Relasi many-to-many Users \& Permissions
    \item \textbf{Role Has Permissions} - Relasi many-to-many Roles \& Permissions
\end{enumerate}

\subsubsection{Entitas Media \& Notifikasi}
\begin{enumerate}
    \item \textbf{Media} - File/gambar polymorphic (users, products, dll)
    \item \textbf{Notifications} - Notifikasi polymorphic untuk users
\end{enumerate}

\subsubsection{Entitas Sistem}
\begin{enumerate}
    \item \textbf{Cache} - Penyimpanan cache
    \item \textbf{Cache Locks} - Lock mechanism untuk cache
    \item \textbf{Jobs} - Queue untuk background jobs
    \item \textbf{Job Batches} - Batch processing jobs
    \item \textbf{Failed Jobs} - Pencatatan job yang gagal
\end{enumerate}

\subsection{Diagram Relasi Entitas}

\begin{verbatim}
┌─────────────────────────────────────────────────────────────────┐
│                         USERS                                   │
│  id (PK) | name | email | password | status | created_at | ... │
└────────────────────────┬──────────────────────────────────────┬─┘
                         │                                      │
         ┌───────────────┼─────────────┬──────────────────────┐ │
         │               │             │                      │ │
    (1:N)│          (1:N)│        (1:N)│                  (1:N)│ │
┌────────▼──────┐ ┌──────▼────────┐ ┌──▼──────────┐ ┌────────▼──────┐
│   SESSIONS    │ │ USER_PROVIDERS│ │ MODEL_HAS_  │ │   NOTIFICATIONS
│               │ │               │ │   ROLES     │ │
│  id (PK)      │ │ id (PK)       │ │             │ │ id (UUID, PK)
│  user_id (FK) │ │ user_id (FK)  │ │ model_id    │ │ notifiable_id
│  ip_address   │ │ provider      │ │ role_id(FK) │ │ notifiable_type
│  user_agent   │ │ provider_id   │ │             │ │ data
│  payload      │ │ avatar        │ │ model_type  │ │ read_at
└───────────────┘ │ created_at    │ └─────────────┘ │ created_at
                  └───────────────┘                  └────────────────┘
                           │
                           │ (1:N)
                  ┌────────▼──────────┐
                  │       MEDIA       │
                  │                   │
                  │ id (PK)           │
                  │ model_id          │
                  │ model_type        │
                  │ collection_name   │
                  │ file_name         │
                  │ mime_type         │
                  │ size              │
                  └───────────────────┘

┌─────────────────────────────────────────┐
│             ROLES                       │
│ id (PK) | name | guard_name | created_at│
└────────────────┬────────────────────────┘
            (1:N)│
                 ├──────────────┬───────────────┐
                 │              │               │
        ┌────────▼────────┐ ┌───▼──────────┐   │
        │ MODEL_HAS_ROLES │ │ ROLE_HAS_    │   │
        │                 │ │ PERMISSIONS  │   │
        │ role_id (FK)    │ │              │   │
        │ model_id        │ │ role_id(FK)  │   │
        │ model_type      │ │ permission_  │   │
        │                 │ │ id(FK)       │   │
        └─────────────────┘ └───┬──────────┘   │
                                │              │
                        ┌───────▼────────┐     │
                        │  PERMISSIONS   │     │
                        │                │     │
                        │ id (PK)        │     │
                        │ name           │◄────┤
                        │ guard_name     │     │
                        │ created_at     │     │
                        └────────────────┘     │
                                               │ (1:N)
                        ┌──────────────────────┘
                        │
                   ┌────▼──────────────────┐
                   │ MODEL_HAS_PERMISSIONS │
                   │                       │
                   │ permission_id (FK)    │
                   │ model_id              │
                   │ model_type            │
                   └───────────────────────┘
\end{verbatim}

\newpage

% ============================================================================
\section{Logical Record Structure (LRS)}
% ============================================================================

Logical Record Structure mendeskripsikan bagaimana data disimpan secara logis dalam bentuk tabel/file.

\subsection{Struktur Tabel Utama}

\subsubsection{USERS}
\begin{tabularx}{\textwidth}{|l|l|l|l|X|}
\hline
\textbf{Atribut} & \textbf{Tipe} & \textbf{Constraint} & \textbf{Indexed} & \textbf{Keterangan} \\
\hline
id & BIGINT & PK, AUTO\_INCREMENT & Yes & Primary Key \\
\hline
name & VARCHAR(255) & NOT NULL & No & Nama lengkap pengguna \\
\hline
email & VARCHAR(255) & UNIQUE, NOT NULL & Yes & Email unik \\
\hline
email\_verified\_at & TIMESTAMP & NULLABLE & No & Waktu verifikasi email \\
\hline
password & VARCHAR(255) & NOT NULL & No & Password terenkripsi \\
\hline
remember\_token & VARCHAR(100) & NULLABLE & No & Token remember me \\
\hline
created\_at & TIMESTAMP & DEFAULT CURRENT\_TIMESTAMP & No & Waktu dibuat \\
\hline
updated\_at & TIMESTAMP & DEFAULT CURRENT\_TIMESTAMP & No & Waktu diupdate \\
\hline
\end{tabularx}

\subsubsection{SESSIONS}
\begin{tabularx}{\textwidth}{|l|l|l|l|X|}
\hline
\textbf{Atribut} & \textbf{Tipe} & \textbf{Constraint} & \textbf{Indexed} & \textbf{Keterangan} \\
\hline
id & VARCHAR(255) & PK & Yes & Session ID \\
\hline
user\_id & BIGINT & FK (Users.id), NULLABLE & Yes & Referensi User \\
\hline
ip\_address & VARCHAR(45) & NULLABLE & No & IP address \\
\hline
user\_agent & TEXT & NULLABLE & No & User agent browser \\
\hline
payload & LONGTEXT & NOT NULL & No & Data session \\
\hline
last\_activity & INT & NOT NULL, INDEXED & Yes & Unix timestamp \\
\hline
\end{tabularx}

\subsubsection{USER\_PROVIDERS}
\begin{tabularx}{\textwidth}{|l|l|l|l|X|}
\hline
\textbf{Atribut} & \textbf{Tipe} & \textbf{Constraint} & \textbf{Indexed} & \textbf{Keterangan} \\
\hline
id & BIGINT & PK, AUTO\_INCREMENT & Yes & Primary Key \\
\hline
user\_id & BIGINT & FK (Users.id), NOT NULL & Yes & Referensi User \\
\hline
provider & VARCHAR(255) & NOT NULL & Yes & Provider (google, github, etc) \\
\hline
provider\_id & VARCHAR(255) & NOT NULL & Yes & ID dari provider \\
\hline
avatar & VARCHAR(255) & NULLABLE & No & Avatar URL \\
\hline
created\_at & TIMESTAMP & DEFAULT CURRENT\_TIMESTAMP & No & Waktu dibuat \\
\hline
updated\_at & TIMESTAMP & DEFAULT CURRENT\_TIMESTAMP & No & Waktu diupdate \\
\hline
\end{tabularx}

\subsubsection{ROLES}
\begin{tabularx}{\textwidth}{|l|l|l|l|X|}
\hline
\textbf{Atribut} & \textbf{Tipe} & \textbf{Constraint} & \textbf{Indexed} & \textbf{Keterangan} \\
\hline
id & BIGINT & PK, AUTO\_INCREMENT & Yes & Primary Key \\
\hline
name & VARCHAR(125) & NOT NULL, UNIQUE & Yes & Nama role \\
\hline
guard\_name & VARCHAR(125) & NOT NULL & Yes & Guard name \\
\hline
created\_at & TIMESTAMP & DEFAULT CURRENT\_TIMESTAMP & No & Waktu dibuat \\
\hline
updated\_at & TIMESTAMP & DEFAULT CURRENT\_TIMESTAMP & No & Waktu diupdate \\
\hline
\end{tabularx}

\subsubsection{PERMISSIONS}
\begin{tabularx}{\textwidth}{|l|l|l|l|X|}
\hline
\textbf{Atribut} & \textbf{Tipe} & \textbf{Constraint} & \textbf{Indexed} & \textbf{Keterangan} \\
\hline
id & BIGINT & PK, AUTO\_INCREMENT & Yes & Primary Key \\
\hline
name & VARCHAR(125) & NOT NULL, UNIQUE & Yes & Nama permission \\
\hline
guard\_name & VARCHAR(125) & NOT NULL & Yes & Guard name \\
\hline
created\_at & TIMESTAMP & DEFAULT CURRENT\_TIMESTAMP & No & Waktu dibuat \\
\hline
updated\_at & TIMESTAMP & DEFAULT CURRENT\_TIMESTAMP & No & Waktu diupdate \\
\hline
\end{tabularx}

\subsubsection{MODEL\_HAS\_ROLES}
\begin{tabularx}{\textwidth}{|l|l|l|l|X|}
\hline
\textbf{Atribut} & \textbf{Tipe} & \textbf{Constraint} & \textbf{Indexed} & \textbf{Keterangan} \\
\hline
role\_id & BIGINT & FK (Roles.id), PK & Yes & Referensi Role \\
\hline
model\_type & VARCHAR(255) & PK & No & Tipe model (App\textbackslash Models\textbackslash User) \\
\hline
model\_id & BIGINT & PK & Yes & ID dari model \\
\hline
\end{tabularx}

\subsubsection{MODEL\_HAS\_PERMISSIONS}
\begin{tabularx}{\textwidth}{|l|l|l|l|X|}
\hline
\textbf{Atribut} & \textbf{Tipe} & \textbf{Constraint} & \textbf{Indexed} & \textbf{Keterangan} \\
\hline
permission\_id & BIGINT & FK (Permissions.id), PK & Yes & Referensi Permission \\
\hline
model\_type & VARCHAR(255) & PK & No & Tipe model \\
\hline
model\_id & BIGINT & PK & Yes & ID dari model \\
\hline
\end{tabularx}

\subsubsection{ROLE\_HAS\_PERMISSIONS}
\begin{tabularx}{\textwidth}{|l|l|l|l|X|}
\hline
\textbf{Atribut} & \textbf{Tipe} & \textbf{Constraint} & \textbf{Indexed} & \textbf{Keterangan} \\
\hline
permission\_id & BIGINT & FK (Permissions.id), PK & Yes & Referensi Permission \\
\hline
role\_id & BIGINT & FK (Roles.id), PK & Yes & Referensi Role \\
\hline
\end{tabularx}

\subsubsection{MEDIA}
\begin{tabularx}{\textwidth}{|l|l|l|l|X|}
\hline
\textbf{Atribut} & \textbf{Tipe} & \textbf{Constraint} & \textbf{Indexed} & \textbf{Keterangan} \\
\hline
id & BIGINT & PK, AUTO\_INCREMENT & Yes & Primary Key \\
\hline
model\_type & VARCHAR(255) & NOT NULL & Yes & Tipe model polymorphic \\
\hline
model\_id & BIGINT & NOT NULL & Yes & ID model polymorphic \\
\hline
uuid & UUID & NULLABLE, UNIQUE & Yes & UUID unik media \\
\hline
collection\_name & VARCHAR(255) & NOT NULL & No & Nama koleksi \\
\hline
name & VARCHAR(255) & NOT NULL & No & Nama file \\
\hline
file\_name & VARCHAR(255) & NOT NULL & No & Nama file lengkap \\
\hline
mime\_type & VARCHAR(255) & NULLABLE & No & MIME type \\
\hline
disk & VARCHAR(255) & NOT NULL & No & Disk storage \\
\hline
size & BIGINT & NOT NULL & No & Ukuran file \\
\hline
created\_at & TIMESTAMP & DEFAULT CURRENT\_TIMESTAMP & No & Waktu dibuat \\
\hline
updated\_at & TIMESTAMP & DEFAULT CURRENT\_TIMESTAMP & No & Waktu diupdate \\
\hline
\end{tabularx}

\subsubsection{NOTIFICATIONS}
\begin{tabularx}{\textwidth}{|l|l|l|l|X|}
\hline
\textbf{Atribut} & \textbf{Tipe} & \textbf{Constraint} & \textbf{Indexed} & \textbf{Keterangan} \\
\hline
id & UUID & PK & Yes & UUID notification \\
\hline
type & VARCHAR(255) & NOT NULL & No & Tipe notifikasi \\
\hline
notifiable\_type & VARCHAR(255) & NOT NULL & Yes & Tipe model polymorphic \\
\hline
notifiable\_id & BIGINT & NOT NULL & Yes & ID model polymorphic \\
\hline
data & TEXT/JSON & NOT NULL & No & Data notifikasi \\
\hline
read\_at & TIMESTAMP & NULLABLE & No & Waktu dibaca \\
\hline
created\_at & TIMESTAMP & DEFAULT CURRENT\_TIMESTAMP & No & Waktu dibuat \\
\hline
updated\_at & TIMESTAMP & DEFAULT CURRENT\_TIMESTAMP & No & Waktu diupdate \\
\hline
\end{tabularx}

\newpage

% ============================================================================
\section{Transformasi ERD ke Logical Relational Schema}
% ============================================================================

\subsection{Aturan Transformasi Umum}

Transformasi dari Entity-Relationship Diagram ke Logical Relational Schema mengikuti aturan-aturan umum berikut:

\begin{enumerate}
    \item \textbf{Setiap Entitas menjadi Relasi (Tabel)}
    \item \textbf{Setiap Atribut menjadi Kolom}
    \item \textbf{Primary Key dipertahankan}
    \item \textbf{Relasi Many-to-Many dikonversi menjadi tabel junction}
    \item \textbf{Foreign Key menunjukkan relasi}
\end{enumerate}

\subsection{Proses Transformasi Digi Herba}

\subsubsection{1. Entitas Menjadi Relasi}

\noindent\textbf{ERD:} Entitas USERS dengan atribut (id, name, email, password, etc)

\noindent\textbf{LRS:}
\begin{verbatim}
USERS (id, name, email, email_verified_at, password, remember_token, 
       created_at, updated_at)
PK: id
\end{verbatim}

\subsubsection{2. Relasi One-to-Many (1:N)}

\noindent\textbf{ERD:} USERS (1) ----- (N) SESSIONS

\noindent\textbf{Transformasi:} Foreign Key dari SESSIONS menunjuk ke USERS

\noindent\textbf{LRS:}
\begin{verbatim}
SESSIONS (id, user_id, ip_address, user_agent, payload, last_activity)
PK: id
FK: user_id REFERENCES USERS(id)
\end{verbatim}

\subsubsection{3. Relasi Many-to-Many (N:M)}

\noindent\textbf{ERD:} USERS (N) ---- (N) ROLES

\noindent\textbf{Transformasi:} Membuat tabel junction (MODEL\_HAS\_ROLES)

\noindent\textbf{LRS:}
\begin{verbatim}
MODEL_HAS_ROLES (role_id, model_type, model_id)
PK: (role_id, model_type, model_id)
FK: role_id REFERENCES ROLES(id)
      model_id REFERENCES USERS(id) [when model_type = 'App\Models\User']
\end{verbatim}

\subsubsection{4. Atribut Multivalued (ROLES memiliki banyak PERMISSIONS)}

\noindent\textbf{ERD:} ROLES (N) ---- (N) PERMISSIONS

\noindent\textbf{Transformasi:} Tabel junction ROLE\_HAS\_PERMISSIONS

\noindent\textbf{LRS:}
\begin{verbatim}
ROLE_HAS_PERMISSIONS (permission_id, role_id)
PK: (permission_id, role_id)
FK: permission_id REFERENCES PERMISSIONS(id)
    role_id REFERENCES ROLES(id)
\end{verbatim}

\subsubsection{5. Relasi Polymorphic (MEDIA ke berbagai model)}

\noindent\textbf{ERD:} MEDIA (1) ----- (N) [USERS, PRODUCTS, etc]

\noindent\textbf{Transformasi:} Dua kolom untuk mengidentifikasi model (model\_type, model\_id)

\noindent\textbf{LRS:}
\begin{verbatim}
MEDIA (id, model_type, model_id, uuid, collection_name, name, 
       file_name, mime_type, disk, size, created_at, updated_at)
PK: id
UNIQUE: uuid
INDEX: (model_type, model_id)

Polymorphic relationship:
  - model_type = 'App\Models\User' → model_id REFERENCES USERS(id)
  - model_type = 'App\Models\Product' → model_id REFERENCES PRODUCTS(id)
\end{verbatim}

\subsection{Foreign Key Constraints}

Semua Foreign Key dalam Logical Relational Schema:

\begin{table}[h]
\centering
\begin{tabular}{|l|l|l|l|}
\hline
\textbf{Tabel} & \textbf{FK Kolom} & \textbf{Referensi} & \textbf{Action} \\
\hline
SESSIONS & user\_id & USERS(id) & ON DELETE CASCADE \\
\hline
USER\_PROVIDERS & user\_id & USERS(id) & ON DELETE CASCADE \\
\hline
MODEL\_HAS\_ROLES & role\_id & ROLES(id) & ON DELETE CASCADE \\
\hline
MODEL\_HAS\_PERMISSIONS & permission\_id & PERMISSIONS(id) & ON DELETE CASCADE \\
\hline
ROLE\_HAS\_PERMISSIONS & permission\_id & PERMISSIONS(id) & ON DELETE CASCADE \\
\hline
ROLE\_HAS\_PERMISSIONS & role\_id & ROLES(id) & ON DELETE CASCADE \\
\hline
\end{tabular}
\end{table}

\subsection{Normalisasi Database}

Database Digi Herba telah dinormalisasi hingga \textbf{3rd Normal Form (3NF)}:

\begin{enumerate}
    \item \textbf{1NF (First Normal Form)}
    \begin{itemize}
        \item Setiap kolom berisi atomic value (tidak ada repeating group)
        \item Semua tabel memiliki primary key
        \item \textcolor{green}{✓ Terpenuhi}
    \end{itemize}
    
    \item \textbf{2NF (Second Normal Form)}
    \begin{itemize}
        \item Sudah 1NF
        \item Setiap non-key attribute bergantung penuh pada primary key
        \item \textcolor{green}{✓ Terpenuhi}
    \end{itemize}
    
    \item \textbf{3NF (Third Normal Form)}
    \begin{itemize}
        \item Sudah 2NF
        \item Tidak ada transitive dependency
        \item Setiap atribut hanya bergantung pada primary key
        \item \textcolor{green}{✓ Terpenuhi}
    \end{itemize}
\end{enumerate}

\newpage

% ============================================================================
\section{Ringkasan Transformasi}
% ============================================================================

\subsection{Mapping Entitas ERD ke Tabel LRS}

\begin{table}[h]
\centering
\begin{tabularx}{\textwidth}{|l|l|c|X|}
\hline
\textbf{ERD Entitas} & \textbf{LRS Tabel} & \textbf{Tipe} & \textbf{Catatan} \\
\hline
USERS & USERS & Entity & Core entity \\
\hline
SESSIONS & SESSIONS & Entity & 1:N with USERS \\
\hline
PASSWORD\_RESET\_TOKENS & PASSWORD\_RESET\_TOKENS & Entity & Transient data \\
\hline
USER\_PROVIDERS & USER\_PROVIDERS & Entity & 1:N with USERS \\
\hline
ROLES & ROLES & Entity & Authorization \\
\hline
PERMISSIONS & PERMISSIONS & Entity & Authorization \\
\hline
USERS $\leftrightarrow$ ROLES & MODEL\_HAS\_ROLES & Junction & N:M relationship \\
\hline
USERS $\leftrightarrow$ PERMISSIONS & MODEL\_HAS\_PERMISSIONS & Junction & N:M relationship \\
\hline
ROLES $\leftrightarrow$ PERMISSIONS & ROLE\_HAS\_PERMISSIONS & Junction & N:M relationship \\
\hline
MEDIA & MEDIA & Entity & Polymorphic (1:N) \\
\hline
NOTIFICATIONS & NOTIFICATIONS & Entity & Polymorphic (1:N) \\
\hline
CACHE & CACHE & System & System table \\
\hline
JOBS & JOBS & System & Queue system \\
\hline
\end{tabularx}
\end{table}

\subsection{Catatan Penting}

\begin{itemize}
    \item \textbf{Polymorphic Relationships:} Menggunakan kombinasi \texttt{model\_type} dan \texttt{model\_id} untuk menunjuk ke berbagai model
    \item \textbf{Cascade Delete:} Foreign keys menggunakan ON DELETE CASCADE untuk menjaga integritas referensial
    \item \textbf{Indexing:} Semua primary key, foreign key, dan kolom yang sering di-query di-index
    \item \textbf{Timestamps:} Setiap tabel memiliki \texttt{created\_at} dan \texttt{updated\_at} untuk audit trail
    \item \textbf{UUID:} Media dan Notifications menggunakan UUID untuk keamanan dan skalabilitas
\end{itemize}

\end{document}
